\chapter{Requirement Analysis and Specifications}

\section{Stakeholders and Roles}
The platform serves multiple stakeholders whose responsibilities define the authorization model:

\subsection{Agency}
Agency users supervise multiple housing units, review owner requests, monitor cross-unit status, and maintain governance controls at the platform level.

\subsection{House Owner / Admin}
House admins manage one household domain: they register devices, supervise residents, review permission requests, activate scenes, and operate safety controls.

\subsection{Resident}
Residents are authenticated users linked to a house code. Their capabilities are intentionally constrained and can be expanded via explicit admin approval.

\subsection{System Operator}
A technical operator (DevOps/maintainer) maintains infrastructure, observability, backups, and updates.

\section{Functional Requirements}

\subsection{Identity and Access}
\begin{itemize}[leftmargin=1.2cm]
  \item Registration and login for resident, admin, and agency flows.
  \item JWT-based API authentication.
  \item Google OAuth login integration.
  \item Email verification and password reset lifecycle.
  \item Route-level and action-level role authorization.
\end{itemize}

\subsection{House and Unit Lifecycle}
\begin{itemize}[leftmargin=1.2cm]
  \item Generation and validation of unit/house codes.
  \item Join existing managed unit or create solo unmanaged unit.
  \item Track claimed/unclaimed unit status.
\end{itemize}

\subsection{Device Management}
\begin{itemize}[leftmargin=1.2cm]
  \item Register IoT devices with metadata (topic, type, room, position).
  \item List and delete devices per house context.
  \item Send command payloads to devices via MQTT topics.
  \item Update online/offline state and last-seen timestamps.
\end{itemize}

\subsection{Permission Workflow}
\begin{itemize}[leftmargin=1.2cm]
  \item Resident submits restricted action request.
  \item Admin reviews and approves/denies.
  \item Approved permissions are merged into resident capabilities.
  \item Notification state supports unread/read transitions.
\end{itemize}

\subsection{Scenarios and Automation}
\begin{itemize}[leftmargin=1.2cm]
  \item Activation of named scenarios (Morning, Dinner, Cinematic, Sleep).
  \item Bulk update of lights and openings according to scene presets.
  \item Synchronization of dashboard state after scenario changes.
\end{itemize}

\subsection{Gas Safety}
\begin{itemize}[leftmargin=1.2cm]
  \item Ingest gas readings from MQTT sensor topics.
  \item Compare readings to configurable threshold.
  \item Trigger emergency mode, alarm command, and valve closure.
  \item Broadcast emergency state to relevant clients in real time.
  \item Allow admin clearing and threshold tuning.
\end{itemize}

\subsection{Realtime Communication}
\begin{itemize}[leftmargin=1.2cm]
  \item Socket.IO authenticated channels for user notifications.
  \item House-room broadcasts for dashboard synchronization.
  \item Agency-global alerts for emergency escalation.
\end{itemize}

\section{Non-Functional Requirements}

\subsection{Performance}
The system should provide low-latency updates for state-changing events, with response times suitable for human-perceived immediacy in dashboard interactions.

\subsection{Scalability}
Architecture should support horizontal evolution in API and socket layers, and future broker scaling strategies.

\subsection{Security}
Critical requirements include password hashing, signed tokens, role isolation, transport security readiness, and controlled data exposure.

\subsection{Maintainability}
Code must be modular, readable, and extensible. Route handlers, models, middlewares, and real-time components should follow clear boundaries.

\subsection{Reliability}
Failure in one communication channel should not corrupt persistent state; safety-critical transitions should remain deterministic and recoverable.

\section{Use Case Inventory}
\begin{longtable}{p{3cm}p{4cm}p{7cm}}
\toprule
Actor & Use Case & Description \\
\midrule
Agency & Review owner requests & Approve/reject account creation requests and issue one-time access code. \\
Agency & Monitor units & View claimed and unclaimed units, statuses, and analytics indicators. \\
Admin & Register device & Add a new smart device and map it to room/topic. \\
Admin & Activate scenario & Apply global scene parameters to eligible devices. \\
Admin & Review request & Approve or deny resident permission request. \\
Admin & Clear emergency & Manually clear emergency mode after hazard assessment. \\
Resident & Request permission & Ask access to restricted actions/rooms. \\
Resident & Control assigned room & Execute allowed actions within authorization scope. \\
System & Process gas reading & Persist reading, evaluate threshold, emit events, publish hardware commands. \\
\bottomrule
\end{longtable}

\section{Acceptance Criteria}
Each major capability has measurable criteria:
\begin{itemize}[leftmargin=1.2cm]
  \item Authentication endpoints return valid token and role payload on success.
  \item Device commands are published to expected MQTT topic format.
  \item Permission approval updates resident permissions and emits real-time update.
  \item Gas reading above threshold transitions to emergency and emits alert events.
  \item Role guards prevent unauthorized access to protected endpoints/pages.
\end{itemize}

\section{Constraints and Assumptions}
\begin{itemize}[leftmargin=1.2cm]
  \item Initial development targets local network and localhost orchestration.
  \item MQTT payload schemas are controlled by platform-owned device firmware simulations.
  \item Users have internet/email access for verification workflows.
  \item Hardware-level commands are simulated in development and can map to physical devices in deployment.
\end{itemize}

\section{Requirement Traceability Overview}
A traceability matrix is maintained in Appendix chapters, mapping each requirement to API routes, model entities, and frontend pages. This ensures auditable engineering alignment from specification to implementation.
